\documentclass{tutodoc}

\usepackage{../preamble.cfg}

\usepackage{../main/ctxt-dsl}
\usepackage{../main/main}


% TESTING LOCAL IMPLEMENTATION %

\usepackage{data}


\begin{document}

\section{Des tableaux de données}

Nous allons voir dans cette section comment saisir de petits tableaux de données de l'un des types suivants.
\begin{itemize}
	\item Tableaux d'images de plusieurs fonctions à une seule variable.

	\item Tableaux d'images d'une seule fonction à deux variables.
\end{itemize}


\begin{tdocnote}
	Tous les tableaux seront fabriqués via \tdocenv{functab} en utilisant un langage spécifique simplifiant la saisie des informations.
	Cet environnement est assez \tdocquote{malin} pour deviner le type de tableau souhaité en fonction des instructions fournies comme nous le constaterons dans cette documentation.
\end{tdocnote}



% --------------- %


\subsection{Tableaux d'images de fonctions à une seule variable}

Les valeurs pivots, qui sont celles dont on veut indiquer les images, s'indiquent de deux façons.
%
\begin{enumerate}
	\item \tdoclatexin{xvals = x1 , x2 , ... , xn} indiquent les valeurs pour la variable nommée $x$ par défaut.

	\item \tdoclatexin{xvals = mavar : x1 , x2 , ... , xn} permet de renommer la variable en $mavar$\,.
\end{enumerate}


Une fois ceci fait, il faut renseigner les différentes images en donnant obligatoirement la formule de l'expression étudiée en passant via \tdoclatexin{imgs = monexpr : i1 , i2 , ... , in}\,.

\medskip

En résumé, tout se passe via \tdoclatexin{xvals = ...} pour les valeurs pivots, et \tdoclatexin{imgs = ...} pour les images.

% --------------- %


\begin{tdocexa}[Une seule fonction avec la variable par défaut]
    \leavevmode

    \tdoclatexinput{examples/1D-data/one-func-no-label.tex}
\end{tdocexa}


\begin{tdocnote}
	Retenir que tout se saisie en mode mathématique.
\end{tdocnote}


\begin{tdocwarn}
	L'utilisation de \tdoclatexin{xvals} doit être faite une fois, et une seule, au tout début du contenu.
\end{tdocwarn}


% --------------- %


\begin{tdocexa}[Deux fonctions pour une variable \tdocquote{maison}]
    \leavevmode

    \tdoclatexinput{examples/1D-data/two-func-label.tex}
\end{tdocexa}


% --------------- %


\begin{tdocexa}[Commentaires à la sauce \LaTeX]
    \leavevmode

    \tdoclatexinput{examples/1D-data/comments.tex}
\end{tdocexa}


% --------------- %

\begin{tdoctip}[Nombres décimaux en version \tdocquote{locale} et \tdocquote{grandes} fractions]
%    \leavevmode
%
    Via les macros \tdocmacro{dfrac} et \tdocmacro{num}
    \footnote{
    	Cette macro ajoute de fins espaces mettant en valeur les groupes de chiffres, tout en gérant l'absence d'espaces autour du séparateur décimal.
    }
    venant des excellents packages \tdocpack{amsmath} et \tdocpack{siuntix} respectivement, il est facile de rédiger des nombres décimaux, et d'obtenir de \tdocquote{grandes} fractions comme le montre l'exemple suivant.

    \tdoclatexinput{examples/1D-data/num-n-dfrac.tex}
\end{tdoctip}

\end{document}
